% Generated draft: requires mathematical and source review.
\documentclass[11pt]{article}
\usepackage[margin=25mm]{geometry}
\usepackage{fontspec}
\setmainfont{Latin Modern Roman}
\usepackage{amsmath,amssymb,mathtools,unicode-math}
\setmathfont{Latin Modern Math}
\usepackage{enumitem,needspace,xurl,hyperref}
\setlength{\parindent}{0pt}
\setlength{\parskip}{5pt}
\setlength{\emergencystretch}{3em}
\newcommand{\N}{\mathbb{N}}
\newcommand{\Z}{\mathbb{Z}}
\newcommand{\Q}{\mathbb{Q}}
\newcommand{\R}{\mathbb{R}}
\newcommand{\C}{\mathbb{C}}
\newcommand{\F}{\mathbb{F}}
\newcommand{\A}{\mathbb{A}}
\newcommand{\PP}{\mathbb{P}}
\newcommand{\E}{\mathbb{E}}
\newcommand{\eps}{\varepsilon}
\newcommand{\dd}{\,\mathrm{d}}
\begin{document}
\section*{UnsolvedMath ID: 30000947}
\textbf{Source problem code:} catalogue entry
\subsection{Definitions and mathematical statement}
Time-domain boundary element methodology (BEM) uses discretizations of boundary integral equations while retaining time as a variable, rather than solving only a fixed-frequency problem. The source asks: Why has time-domain BEM not achieved the operational deployment expected a decade ago, and what mathematical and computational developments are needed to overcome the stability, precision, and cost barriers for time-domain integral equations? This preserves both the explanatory deployment question and the request for developments addressing all three barriers. The historical premise and a decade ago belong to the supplied report's context; they are not asserted as an independently verified description of deployment in 2026. No particular underlying PDE, boundary geometry, causal kernel, discretization, stability norm, error tolerance, complexity target, or operational benchmark is given. Consequently the text specifies a numerical-analysis and engineering research direction rather than a single quantified theorem or algorithmic specification. The 2026-08-21 literature triage labels the record open and says the report identifies stability and implementation barriers rather than a theorem resolving the deployment question.
\subsection{Short English statement}
Explain the report's shortfall in operational use of time-domain BEM and identify developments addressing stability, precision, and computational cost. The source supplies a deployment research question, with no particular equation or quantitative performance target.
\subsection{Sources}
\begin{itemize}\item Local UnsolvedMath registry, permanent ID 30000947, source problem catalogue entry, \emph{Deployment Barriers for Time-Domain Boundary Element Methods}, proposed year 2008; no individual proposer supplied.\item Supplied citation: \emph{Analysis of Boundary Element Methods} (2009), Oberwolfach Reports, DOI 10.4171/owr/2008/19, pp.\ 953--1030: \url{https://doi.org/10.4171/owr/2008/19}. Supplied literature label: Oberwolfach Report 19/2008.\end{itemize}
\end{document}
